\section{סיכום אינטגרטיבי: המשמעות הניהולית והטכנולוגית}

פרק זה מחבר את הנושאים הטכניים למסגרת החלטות ארגונית: מה משתנה כשארגון
מאמץ סוכני \eng{AI} לא רק כצ'אט, אלא כשכבת ביצוע?

\subsection{מעבר מ-GUI ל-API}

המגמה המרכזית מההרצאה: תוכנות יחשפו יכולות בטקסט ו-\eng{API} במקום
להסתמך על \eng{GUI} בלבד. ארגונית:
\begin{itemize}[rightmargin=2em]
  \item צוותי פיתוח --- עדיפות ל-\eng{API} ו-\eng{MCP} על פני אוטומציית מסך.
  \item צוותי תפעול --- הגדרת \eng{Skills} ותבניות מסמך (\eng{CLS}).
  \item הנהלה --- ציפייה לרפטביליות ולא ל``קסם'' חד-פעמי של סוכן.
\end{itemize}

\subsection{שלוש שכבות בארגון}

ניתן למפות את מודל השכבות לתפקידים:
\begin{description}[rightmargin=2em, style=nextline]
  \item[Software] מערכות ליבה, \eng{ERP}, מאגרי נתונים, קומפיילרים.
  \item[Skill] מפרטי עבודה: איך לכתוב דוח, איך לבדוק טבלה, איך לסכם פגישה.
  \item[Agent] נקודת כניסה למשתמש --- מממש מטרות בשפה טבעית.
\end{description}

\subsection{מודל עלויות}

\begin{itemize}[rightmargin=2em]
  \item טוקנים --- חלון הקשר, צילומי מסך, היסטוריה ארוכה.
  \item זמן --- סוכן יחיד לעומת מקביליות עם אורכוסטרטור.
  \item תחזוקה --- \eng{Skills} ו-\eng{CLS} מקטינים עלות תיקון חוזר.
\end{itemize}

\subsection{רפטביליות כיתרון תחרותי}

ארגון שבנה:
\begin{itemize}[rightmargin=2em]
  \item ספריית \eng{Skills} מאושרת,
  \item תבניות \eng{CLS}/\eng{PDF} אחידות,
  \item מערכת \eng{QA} היררכית,
\end{itemize}
יכול לייצר עשרות דוחות, הצעות מחיר או מסמכי רגולציה באותה איכות ---
במרחק פרומפט. מתחרה שעדיין עורך \eng{Word} ידנית נשאר מאחור.

\subsection{סיכונים וניהול סיכונים}

\begin{itemize}[rightmargin=2em]
  \item הזיות מודל --- במיוחד ב-\eng{CTI} ובמשפט; נדרש אימות אנושי.
  \item דליפת מידע --- פרומפטים עם נתונים רגישים.
  \item תלות בספק --- מודל, \eng{MCP}, ענן.
  \item ריבוי סוכנים ללא פיקוח --- לולאות ועלויות.
\end{itemize}

\agenttable{
\caption{מסגרת החלטה ארגונית לסוכן \eng{AI}}
\label{tab:org-framework}
\begin{tabular}{@{}p{3cm}p{4.5cm}p{4.5cm}@{}}
\toprule
\textbf{שאלה} & \textbf{אם כן} & \textbf{אם לא} \\
\midrule
האם חוזר על עצמו? & \eng{CLS} + \eng{Skills} & חד-פעמי ב-\eng{GUI} \\
האם דורש ידע פנימי? & \eng{RAG} מאובטח & פרומפט כללי \\
האם מחובר לתוכנה? & \eng{API}/\eng{MCP} & סוכן דפדפן \\
האם דורש מהירות? & ריבוי סוכנים + אורכוסטרטור & סוכן יחיד \\
\bottomrule
\end{tabular}
}

\subsection{קשר לתפקיד מומחה AI בארגון}

המרצה ציפה שבוגרי הקורס יהיו ``מובילי הטמעת \eng{AI}'' --- לא רק משתמשי
צ'אט, אלא מגדירי תבניות, \eng{Skills}, תהליכי \eng{QA} וארכיטקטורת
אינטגרציה. מטלה 04 מאמנת בדיוק מיומנות זו: הגדרת \eng{PRD}, בניית \eng{CLS},
והפקת תוצר מקצועי.

\begin{insightbox}
המעבר מטכנולוגיה לערך עסקי: סוכן שמייצר \eng{PDF} אחיד ב-30 שניות
מחליף שעות עיצוב --- אם התשתית (\eng{CLS}, \eng{QA}) כבר קיימת.
\end{insightbox}
